Large language model (LLM) based autonomous agents have advanced rapidly in recent years, demonstrating remarkable capabilities in code generation, web browsing, and complex task automation~\cite{react,cot,gpt4}.
Systems such as Claude Code~\cite{claudecode}, Codex~\cite{codex}, and OpenHands~\cite{openhands} can now interpret high-level instructions and execute multi-step plans that span dozens of tool invocations.
Yet the vast majority of these systems are architected for \emph{single-session} task completion: an agent receives a request, carries it out, and terminates without retaining any knowledge for future sessions.
This design assumption breaks down when agents are expected to operate autonomously over extended time horizons---managing personal workflows, monitoring evolving environments, or accumulating domain expertise across hundreds of interactions.
In such long-term settings, the context window becomes the fundamental bottleneck rather than model capability itself~\cite{anthropic_context}.
Tool definitions, historical memory, retrieved documents, and raw web content all compete for the same finite token budget, continuously expanding the prompt while squeezing the space available for effective reasoning.
Empirical evidence consistently shows that, given the same informational requirements, longer prompts degrade model performance: attention becomes diluted, irrelevant tokens introduce noise, and critical instructions are more likely to be overlooked.
This tension between the growing information demands of long-term operation and the fixed capacity of the context window constitutes the central challenge that motivates our work.

We identify four concrete dimensions along which context information density is eroded in existing agent systems, each representing a distinct source of what we term \emph{context pollution}.
First, \textbf{memory retrieval}: as an agent accumulates knowledge across sessions, its memory store grows without bound, yet each decision typically requires only a small, precise subset of that knowledge---retrieving too much wastes context, while retrieving too little causes the agent to forget critical facts.
Second, \textbf{experience reuse}: when an agent encounters a task similar to one it has solved before, the ideal behavior is to leverage prior experience; however, raw reasoning traces are verbose, filled with exploratory dead ends and redundant deliberation that are costly to re-inject into the prompt.
Third, \textbf{tool overhead}: standard agent frameworks inject the full schema of every available tool into the system prompt on every LLM call, imposing a persistent, per-turn token cost that scales linearly with the size of the tool library~\cite{react}.
Fourth, \textbf{web interaction}: when agents browse the web, raw HTML pages contain navigation bars, advertisements, scripts, and boilerplate that vastly outnumber the task-relevant content, yet conventional approaches pass this entire payload to the model.
Existing systems have addressed these challenges in isolation---MemGPT~\cite{memgpt} proposes OS-inspired virtual memory but lacks mechanisms for self-evolution, Voyager~\cite{voyager} introduces a skill library but is confined to Minecraft, and multi-agent architectures such as Manus~\cite{manus} distribute complexity across specialized sub-agents without fundamentally reducing per-agent context consumption.
No prior system has addressed all four dimensions under a unified design principle.

In this paper, we present \textbf{GenericAgent}, a general-purpose LLM agent system designed from the ground up around the principle of \emph{context information density maximization}: within the finite context window available to each LLM call, the system should pack the highest-density, most task-relevant information so as to maximize the probability of task success while minimizing total prompt cost.
The corollary of this principle is equally important: when the information required for a decision is already complete, a shorter prompt will yield better performance than a longer one, because every superfluous token is a source of distraction rather than signal.
GenericAgent instantiates this principle through four tightly integrated subsystems---hierarchical memory, reflection-driven self-evolution, minimal atomic tools, and context-compressed browser control---each targeting one of the four information density challenges identified above.
All subsystems are unified under a single \emph{Agent Loop} that drives three complementary operating modes: \textbf{Interactive} mode for direct user commands, \textbf{Task} mode for sub-agent coordination via structured file I/O, and \textbf{Reflect} mode for event- or cron-triggered autonomous operation.
The entire system is implemented in approximately 3,000 lines of core code, supports multiple LLM backends including Claude, OpenAI, and Gemini, and is publicly available at \url{https://github.com/lsdefine/GenericAgent}.

Our main contributions are as follows:
\begin{itemize}[leftmargin=*]
    \item A \textbf{hierarchical memory system} (L0--L3) with taxonomy-based indexing and action-verified updates, ensuring that only the minimal, validated information needed for each decision is retrieved into the context window.
    \item A \textbf{reflection-driven self-evolution} mechanism that progressively compresses raw experience into standard operating procedures, then into executable code, and ultimately into composable skills---each stage reducing the context cost of reusing past knowledge by an order of magnitude.
    \item A set of \textbf{seven minimal atomic tools} whose compact definitions occupy a fraction of the prompt budget consumed by conventional tool libraries, yet achieve full environmental control through principled composition~\cite{reflexion}.
    \item A \textbf{context-compressed browser control} scheme based on simplified HTML rendering and incremental DOM monitoring, which strips web pages down to their task-relevant semantic content before injection into the context window.
\end{itemize}
